\documentclass[11pt]{article}
\usepackage{amsmath,amssymb,amsthm,hyperref}
\newtheorem{corollary}{Corollary}
\begin{document}
\title{The exact Gaussian dimension--degree boundary: a sourced corollary}
\author{}\date{17 November 2026}\maketitle

Let $G(n,d)$, for integers $n\geq1$ and $d\geq0$, mean that every
complex polynomial $P$ of total degree at most $d$ in $n$ independent
standard real Gaussian variables, satisfying $\mathbb E[P^m]=0$ for
every $m\geq1$, satisfies $\mathbb E[QP^m]=0$ for all sufficiently
large $m$, for every fixed polynomial $Q$.

\begin{corollary}
With the exact scopes of the inputs below,
\[
 G(n,d)\quad\Longleftrightarrow\quad
 n\leq2\quad\hbox{or}\quad d\leq2\quad\hbox{or}\quad(n,d)=(3,3).
\]
For the ternary-cubic positive case the sufficient strict threshold is
$m>3\deg Q$, and its coefficient three is sharp.
\end{corollary}

\paragraph{Inputs and credit.}
Wilson's public preprint arXiv:2607.23887v1, Theorem~1, proves the
unrestricted two-variable assertion, with threshold $m>\deg Q$.
Liu's public preprint, Zenodo~21688251 version~3, gives the quadratic
theorem in every dimension. The retained internal theorem
\texttt{G3CUBICALL20261115}, with canonical source
\texttt{TERNARY\_CUBIC\_GMC\_CLOSURE.md}, supplies the ternary-cubic
case and its sharp mixed threshold. Long's public preprint
arXiv:2607.18186v1, Proposition~4.1 and Theorem~5.1, supplies failures
of degree three in four variables and degree four in three variables.
The user's 17 November observation identified the exact combined boundary.

\begin{proof}
The two-variable result also covers one variable by ignoring the second
Gaussian coordinate. The quadratic theorem and the ternary-cubic theorem
give the other stated positive ranges.

If none of these ranges applies, then $n\geq3$ and $d\geq3$.
For $n=3$ we must have $d\geq4$, so Long's ternary quartic applies.
For $n\geq4$, Long's four-variable cubic applies. Ignore additional
independent coordinates in both the phase and its witnessing multiplier;
this preserves all their pure and mixed moments. Increasing a degree
bound preserves any counterexample already below it. This exhausts
the complement of the positive ranges.

The sharp threshold is part of the retained cubic theorem, and can be
checked separately from its qualitative reduction. In circular Gaussian
coordinates, $P=W+ZT-Z^3/2$ has
\[
 \mathbb E[e^{tP+\beta W}]
   =\exp[-\beta t(t+\beta)^2/2].
\]
Thus $P$ is pure and
$\mathbb E[W^D P^{3D}]=(3D)!(-1/2)^D\ne0$ for every $D\geq1$.
No coefficient smaller than three can therefore work uniformly.
\end{proof}

\paragraph{Assurance and manuscript scope.}
This is a derived corollary of the credited external inputs and the
existing internal ternary theorem, not a new independent verification
of those proofs. All independent replay, formal verification and external
review flags remain false; priority has not been assessed. The theorem
classifies universal dimension--degree ranges, not every individual safe
phase. The working manuscript in
\texttt{papers/gaussian-boundary/manuscript.tex} has two explicitly
unfinished standalone cubic proof consolidations. The long-term writing
objective in \texttt{papers/gaussian-boundary/paper\_goal.json} remains
active until those proofs and the readability/source audits are complete.
All earlier proof and checker packets are unchanged.
\end{document}
